\documentclass[12pt]{article}

\usepackage{amsmath,amssymb}
\usepackage{nomencl}
\usepackage{etoolbox}

\makenomenclature

%------------------------------------------------
% Group headings
%------------------------------------------------
\renewcommand{\nomgroup}[1]{%
	\item[\bfseries
	\ifstrequal{#1}{P}{Physics Constants}{%
		\ifstrequal{#1}{N}{Number Sets}{%
			\ifstrequal{#1}{M}{Mathematical Operators}{%
				\ifstrequal{#1}{A}{Astronomical Symbols}{%
					\ifstrequal{#1}{U}{Units of Measurement}{%
						\ifstrequal{#1}{O}{Other Symbols}{}}}}}}%
	]
}

% Title
\renewcommand{\nomname}{List of Symbols and Constants}

\begin{document}
	
	\tableofcontents
	
	\section{Introduction}
	
	This document demonstrates a grouped nomenclature.
	
	%------------------------------------------------
	% Physics Constants
	%------------------------------------------------
	\nomenclature[P]{$c$}{Speed of light in vacuum}
	\nomenclature[P]{$h$}{Planck constant}
	\nomenclature[P]{$\hbar$}{Reduced Planck constant}
	\nomenclature[P]{$G$}{Gravitational constant}
	\nomenclature[P]{$\mu_0$}{Permeability of free space}
	\nomenclature[P]{$\epsilon_0$}{Permittivity of free space}
	
	%------------------------------------------------
	% Number Sets
	%------------------------------------------------
	\nomenclature[N]{$\mathbb{N}$}{Natural numbers}
	\nomenclature[N]{$\mathbb{Z}$}{Integers}
	\nomenclature[N]{$\mathbb{Q}$}{Rational numbers}
	\nomenclature[N]{$\mathbb{R}$}{Real numbers}
	\nomenclature[N]{$\mathbb{C}$}{Complex numbers}
	\nomenclature[N]{$\mathbb{H}$}{Quaternions}
	
	%------------------------------------------------
	% Mathematical Operators
	%------------------------------------------------
	\nomenclature[M]{$\nabla$}{Gradient operator}
	\nomenclature[M]{$\Delta$}{Laplacian operator}
	\nomenclature[M]{$\partial$}{Partial derivative}
	\nomenclature[M]{$\sum$}{Summation operator}
	\nomenclature[M]{$\prod$}{Product operator}
	\nomenclature[M]{$\int$}{Integral operator}
	
	%------------------------------------------------
	% Astronomical Symbols
	%------------------------------------------------
	\nomenclature[A]{$M_{\odot}$}{Solar mass}
	\nomenclature[A]{$R_{\odot}$}{Solar radius}
	\nomenclature[A]{$L_{\odot}$}{Solar luminosity}
	\nomenclature[A]{$M_{\oplus}$}{Earth mass}
	\nomenclature[A]{$R_{\oplus}$}{Earth radius}
	\nomenclature[A]{AU}{Astronomical Unit}
	
	%------------------------------------------------
	% Units
	%------------------------------------------------
	\nomenclature[U]{m}{metre}
	\nomenclature[U]{kg}{kilogram}
	\nomenclature[U]{s}{second}
	\nomenclature[U]{A}{ampere}
	\nomenclature[U]{K}{kelvin}
	\nomenclature[U]{mol}{mole}
	
	%------------------------------------------------
	% Other Symbols
	%------------------------------------------------
	\nomenclature[O]{$V$}{Volume}
	\nomenclature[O]{$\rho$}{Density}
	\nomenclature[O]{$\sigma$}{Stefan--Boltzmann constant}
	\nomenclature[O]{$\theta$}{Angular displacement}
	\nomenclature[O]{$\omega$}{Angular frequency}
	\nomenclature[O]{$\lambda$}{Wavelength}
	
	\section{Applications}
	
	The speed of light $c$ appears in Einstein's equation
	
	\[
	E = mc^2.
	\]
	
	The Planck constant $h$ is fundamental in quantum mechanics.
	
	The set of real numbers is denoted by $\mathbb{R}$.
	
	\clearpage
	
	\printnomenclature
	
\end{document}